% Questi appunti affrontano argomenti  della matematica discreta. Il contenuto spazia dalla logica   fino a strumenti sofisticati per il calcolo combinatorio e il principio di induzione, fornendo un quadro completo  delle idee che stanno alla base di questo ambito della matematica.

% \section{Considerazioni Finali}
% Questi appunti sono stati elaborati con l'obiettivo di fornire uno strumento integrativo allo studio, utile sia come guida teorica che come riferimento per esercizi pratici. Tuttavia, essi non intendono sostituire testi accademici o le lezioni frontali, ma piuttosto integrarli, rendendo i concetti più accessibili e organizzati.



% \begin{thebibliography}{99}

% \bibitem{funzione di Cantor a due variabil} 
% Wikipedia. \emph{Funzione di Cantor a due variabi}.  
% Disponibile su: \url{https://en.wikipedia.org/wiki/Pairing_function}

% \bibitem{Lancia}
% Giuseppe Lancia. \emph{Matematica Discreta}.  2022. ISBN: 978-8847024919.

% \bibitem{Epp}
% Susanna S. Epp. \emph{Discrete Mathematics with Applications}, 4th Edition. Cengage Learning, 2010. ISBN: 978-0495391326.

% \bibitem{Rosen}
% Kenneth H. Rosen. \emph{Discrete Mathematics and Its Applications}, 7th Edition. McGraw-Hill Education, 2019. ISBN: 978-0073383095.

% \bibitem{AppuntiSlide}
% Sabina Rossi. \emph{Appunti e Slide delle lezioni di strutture discrete }, Anno Accademico 2024-2025.

% \end{thebibliography}